% !TeX root = ../../main.tex
\subsection{Phân tích yêu cầu}
\subsubsection{Yêu cầu chức năng}

Yêu cầu chức năng được xây dựng từ các quyết định ở Bảng~\ref{tab:quyetdinh}: mỗi nhóm người dùng cần xem được những thông tin nào, ở mức chi tiết nào.

\paragraph{Đối với ban điều hành sàn}
\begin{itemize}
\item Xem tổng quan doanh thu, số đơn hàng, giá trị đơn hàng trung bình và điểm đánh giá trung bình của toàn sàn.
\item So sánh các chỉ số theo tháng, quý, năm và xem xu hướng tăng trưởng.
\item Phân tích doanh thu theo bang, thành phố và ngành hàng, đi từ mức tổng hợp xuống mức chi tiết.
\end{itemize}

\paragraph{Đối với quản lý vận hành và giao hàng}
\begin{itemize}
\item Theo dõi thời gian giao hàng trung bình và tỷ lệ giao trễ theo thời gian, bang và ngành hàng.
\item Xem tỷ lệ trễ theo tuyến giao hàng, tức từ bang của người bán đến bang của khách hàng.
\item Tách được thời gian của từng công đoạn: duyệt đơn, chuẩn bị hàng và vận chuyển.
\item So sánh ngày giao thực tế với ngày giao dự kiến.
\end{itemize}

\paragraph{Đối với quản lý ngành hàng và người bán}
\begin{itemize}
\item Xếp hạng người bán và ngành hàng theo doanh thu, số đơn hàng, tỷ lệ giao trễ và điểm đánh giá, có lọc theo vùng, ngành hàng và khoảng thời gian.
\item Tìm ra những người bán có doanh thu cao nhưng chất lượng dịch vụ thấp.
\item So sánh nơi người bán tập trung với nơi khách hàng tập trung.
\end{itemize}

\paragraph{Đối với bộ phận chăm sóc khách hàng}
\begin{itemize}
\item Theo dõi điểm đánh giá trung bình và phân bố điểm từ 1 đến 5 sao.
\item Xem tỷ lệ đánh giá thấp theo ngành hàng, người bán và khu vực.
\item So sánh điểm đánh giá giữa đơn giao đúng hạn và đơn giao trễ.
\end{itemize}

\paragraph{Đối với quy trình ETL và kho dữ liệu}
Về phía hệ thống, quy trình ETL phải đọc được 9 tập tin CSV nguồn, làm sạch và tích hợp dữ liệu theo các quy tắc xử lý cho từng vấn đề trong Bảng~\ref{tab:chatluong}, và tính sẵn các trường dẫn xuất cần cho các chỉ số ở Bảng~\ref{tab:kpi}. Dữ liệu được lưu theo mô hình đa chiều để thực hiện được các phép toán OLAP \texttt{slice}, \texttt{dice}, \texttt{drill-down}, \texttt{roll-up} và \texttt{pivot}. Mỗi lần chạy ETL phải để lại nhật ký, kể cả số bản ghi bị loại hoặc bị sửa. Kho dữ liệu cũng là nguồn dữ liệu duy nhất cho dashboard và thành phần DSS.

\subsubsection{Yêu cầu phi chức năng}

Các yêu cầu phi chức năng dưới đây được đặt ở mức kiểm chứng được; kết quả đo thực tế được trình bày ở phần hiện thực (mục~\ref{sec:hienthuc}).

\paragraph{Performance -- Hiệu năng}
\begin{itemize}
\item Các truy vấn tổng hợp phục vụ dashboard trả kết quả trong tối đa 5 giây.
\item Một lần chạy ETL đầy đủ với dữ liệu hiện có hoàn thành trong tối đa 30 phút.
\end{itemize}

\paragraph{Data Quality -- Chất lượng dữ liệu}
\begin{itemize}
\item Số liệu tổng trong kho khớp với dữ liệu nguồn, sau khi trừ các bản ghi bị loại có chủ đích.
\item Mọi bản ghi trong bảng fact đều tham chiếu đến một bản ghi hợp lệ trong bảng chiều (dimension table).
\item Bản ghi bị loại hoặc bị sửa đều được ghi lại và truy vết được.
\end{itemize}

\paragraph{Scalability -- Khả năng mở rộng}
\begin{itemize}
\item Có thể thêm nguồn dữ liệu, chiều phân tích (dimensions) hoặc thước đo (measures) mới mà không phải thay đổi toàn bộ cấu trúc.
\item Thiết kế không phụ thuộc vào quy mô dữ liệu hiện tại; khi dữ liệu tăng có thể bổ sung chỉ mục, phân vùng hoặc bảng tổng hợp lưu sẵn.
\end{itemize}

\paragraph{Reliability -- Tính tin cậy}
\begin{itemize}
\item Chạy lại quy trình ETL không làm trùng lặp dữ liệu.
\item Khi một bước ETL lỗi, hệ thống dừng lại, ghi nhận lỗi và không để lại dữ liệu nạp dang dở trong kho.
\end{itemize}

\paragraph{Security -- Tính bảo mật}
\begin{itemize}
\item Người dùng phân tích chỉ có quyền đọc; quyền ghi chỉ dành cho quy trình ETL.
\item Thông tin kết nối cơ sở dữ liệu không được lưu trong mã nguồn.
\end{itemize}

\paragraph{Usability -- Tính khả dụng}
\begin{itemize}
\item Nhà quản lý không cần biết SQL vẫn lọc và xem được dữ liệu trên dashboard.
\item Tên chiều và thước đo phản ánh đúng ý nghĩa nghiệp vụ.
\end{itemize}

\paragraph{Maintainability -- Tính duy trì}
\begin{itemize}
\item Có tài liệu mô tả cấu trúc kho, nguồn gốc dữ liệu và quy tắc chuyển đổi (metadata).
\item Mã nguồn ETL được tách thành các bước độc lập, dễ kiểm tra và chỉnh sửa.
\end{itemize}
